\subsection{Characterizing Gene-Trait Relevance Across the Phenome}

% Narrative arc:
% 1. How many genes does PIGEAN implicate? (scope)
% 2. Among those, what is known vs new? (mechanism vs discovery)
% 3. Aggregating across traits: are pleiotropic genes under constraint? (pleiotropy-constraint)
% 4. Which genes are most pleiotropic, and what do they tell us? (characterization)
% 5. Does this generalize across disease domains? (trait stratification)

Having validated that indirect support predicts direct genetic evidence and prioritizes successful drug targets, we next asked what biological properties characterize the gene-trait relevance landscape that PIGEAN produces. We assessed the scope of evidence across 18,325 genes and 3,026 traits by taking the maximum evidence score per gene across all traits. Among common variant traits, 50\% of genes (9,158) had strong direct support (log Bayes factor $> 3$) for at least one trait, and 18\% (3,331) had strong indirect support (prior $> 3$). Combining both evidence types raised coverage to 50\% of genes (9,246 above a combined threshold of 3), demonstrating that indirect evidence identifies genes beyond those detectable by direct GWAS signal alone. For rare disease traits, where direct variant-level evidence is not available, 25\% of genes (4,196) had indirect support above threshold, indicating that pathway information provides graded discrimination among candidate disease genes even in the absence of association data.

Given this broad coverage, we next asked how much of PIGEAN's evidence reflects known biology versus potential new discoveries. We classified common variant gene--trait pairs by whether they had direct support (log Bayes factor $> 3$), indirect support (prior $> 3$), both, or neither. Across 4.6 million gene--trait pairs, 0.13\% (5,971) were classified as understood mechanisms (both direct and indirect support), 1.02\% (46,409) as novel discoveries (direct support without indirect corroboration), and 0.43\% (19,453) as indirect-only associations. Notably, among the 52,380 pairs with direct support, 89\% lacked corresponding pathway annotation, suggesting that current gene set databases capture only a fraction of the biology underlying GWAS signals. The highest-confidence understood mechanisms for portal traits recapitulated canonical biology: \textit{ANK1}--immature reticulocyte fraction (direct 5.01, indirect 7.91), \textit{INS}--type 2 diabetes (5.11, 7.23), \textit{KLF1}--reticulocyte traits (5.11, 7.08), \textit{LDLR}--dyslipidemia (5.10, 6.72), and \textit{APOE}--myocardial infarction (4.81, 6.73)---genes residing in well-curated pathways (erythropoiesis, insulin signaling, cholesterol metabolism) that corroborate the direct GWAS signal.

Novel discoveries, by contrast, comprised gene--trait pairs with strong direct evidence but low or absent indirect support, indicating associations not captured by current pathway databases. Applying a stringent indirect threshold (prior $< 0.5$) identified 9,738 such pairs among portal traits alone. Several represent biologically compelling associations where the gene's function is directly relevant to the phenotype yet absent from curated gene sets: \textit{GOT1}--AST/ALT ratio (direct 4.75, indirect 0.42), where the gene encodes the enzyme that \emph{is} serum AST; \textit{ELN}--aortic root velocity (direct 4.76, indirect 0.38), where elastin is the primary structural protein of arterial walls; \textit{ATOH7}--optic cup area (direct 4.71, indirect 0.04), a transcription factor essential for retinal ganglion cell development; and \textit{GHSR}--waist circumference (direct 4.73, indirect 0.12), encoding the ghrelin receptor that regulates appetite and energy balance. Pairs with negative indirect scores represent even more striking gaps in pathway databases: \textit{ZHX3}--fasting glucose (direct 4.64, indirect $-$0.16), \textit{GREB1}--uterine fibroids (direct 4.54, indirect $-$0.01), and \textit{KCNRG}--primary biliary cholangitis (direct 4.64, indirect $-$0.22). At the gene level, the most broadly novel loci included \textit{TRIB1} (37 portal trait associations with indirect $< 0.5$), \textit{CENPW} (36), and \textit{PPP1R3B} (34)---pleiotropic GWAS loci whose breadth of association across cardiometabolic phenotypes substantially outstrips their current pathway annotation.

The observation that many genes contribute to multiple traits raised a natural question: are pleiotropic genes under stronger evolutionary constraint? To test this, we computed PIGEAN-derived pleiotropy scores across all 3,026 traits and 18,325 genes, applying inverse-correlation weighting to down-weight redundant phenotypes (reducing the effective number of independent traits to approximately 130 for combined evidence, 311 for direct, and 61 for indirect). We then regressed these scores against gnomAD v4 constraint metrics ($N \approx 15{,}800$ genes with constraint data). Using combined evidence, we observed a significant negative correlation between LOEUF and pleiotropy (Spearman $\rho = -0.261$, $P < 10^{-244}$; Extended Data Table~\ref{tab:pleiotropy_constraint}), indicating that genes more intolerant to loss-of-function variants exhibit greater pleiotropy. Loss-of-function z-score showed a concordant positive association ($\rho = 0.250$, $P < 10^{-224}$), missense constraint was more weakly associated (mis\_z $\rho = 0.171$, $P < 10^{-104}$), and synonymous z-score---included as a negative control---showed no meaningful association ($\rho = 0.009$, $P = 0.23$), confirming that the relationship reflects genuine selective pressure rather than technical confounders. Categorizing genes by LOEUF threshold confirmed this: constrained genes (LOEUF $< 0.35$; $n = 1{,}294$) had nearly twice the mean pleiotropy of tolerant genes (LOEUF $> 1.0$; $n = 7{,}232$; $\mu = 217$ vs.\ $114$; Welch's $t = 15.6$, $P < 10^{-50}$; Mann--Whitney $P < 10^{-97}$), and this difference was significant across all score types and trait categories (all $P < 10^{-12}$; Extended Data Figure~\ref{fig:pleiotropy-violin}). Both direct and indirect evidence independently recapitulated this association (direct LOEUF $\rho = -0.252$; indirect $\rho = -0.214$), consistent with constrained genes accumulating pleiotropic evidence through complementary mechanisms: direct genetic associations at the locus level and participation in trait-associated biological pathways. The association was consistent across portal traits (LOEUF $\rho = -0.247$), GWAS catalog traits ($\rho = -0.208$), and rare disease traits ($\rho = -0.206$), demonstrating that the constraint--pleiotropy relationship generalizes across disease domains (Extended Data Table~\ref{tab:pleiotropy_constraint}).

We next examined which specific genes drive this pattern. Across all traits using combined evidence, the most pleiotropic genes were \textit{TP53} (score 2{,}291), \textit{APOE} (2{,}060), \textit{CTNNB1} (2{,}001), \textit{PTEN} (1{,}922), and \textit{SHH} (1{,}776)---all highly constrained master regulators that participate in diverse biological processes. Decomposing this signal revealed distinct modes of pleiotropy. Direct evidence was dominated by HLA class II genes (\textit{HLA-DRB1}, \textit{HLA-DQB1}, \textit{HLA-B}) and broadly associated GWAS loci (\textit{SH2B3}, \textit{GCKR}, \textit{TERT}), reflecting genes with strong, detectable variant effects across many traits. Indirect evidence, by contrast, was dominated by developmental signaling hubs---\textit{TP53}, \textit{CTNNB1}, \textit{SHH}, \textit{BMP4}, \textit{GLI3}, \textit{NOTCH1}---genes whose pleiotropy derives from central positions in biological pathway networks rather than from individually significant GWAS associations. This decomposition illustrates that the direct score identifies genomic regions of broad effect, while the indirect score identifies genes that are biologically central even when their individual variant effects are modest.

Finally, stratification by trait category revealed domain-specific patterns of pleiotropy. Among GWAS catalog traits, the most pleiotropic genes were dominated by lipid metabolism (\textit{APOE}, \textit{APOB}, \textit{CETP}, \textit{LDLR}, \textit{LPL}, \textit{APOA1}), consistent with the well-powered lipid GWAS that contribute disproportionately to this trait set. Portal traits highlighted metabolic disease hubs (\textit{APOE}, \textit{PPARG}, \textit{INS}, \textit{LEP}, \textit{LEPR}, \textit{IRS1}) alongside developmental regulators (\textit{PTPN11}, \textit{SHH}, \textit{FGFR2}). Most strikingly, rare disease traits were enriched for ciliopathy genes (\textit{CEP290}, \textit{MKS1}, \textit{IFT172}, \textit{RPGRIP1L}, \textit{BBS4}) and RASopathy genes (\textit{KRAS}, \textit{BRAF}, \textit{PTPN11}, \textit{NF1})---Mendelian disease genes whose mutations cause multi-system syndromes affecting diverse organ systems. This represents genuine biological pleiotropy: ciliary signaling is required for the development of kidneys, eyes, brain, and limbs, and mutations in ciliary genes accordingly produce syndromic phenotypes spanning these systems. The convergence of PIGEAN's data-driven pleiotropy scores with known syndromic biology provides strong biological validation of the approach.

% TODO: Add paragraph on genes that are both highly pleiotropic AND highly
% constrained - what is unique about this intersection? Are these the
% "core disease genes"? How do they compare to less constrained pleiotropic genes?

% FIGURE SUGGESTION: A scatter plot of pleiotropy score vs LOEUF colored by
% gene category (e.g., known drug targets, essential genes, transcription factors)
% would visualize the constraint-pleiotropy relationship and highlight
% biologically interesting outliers.
